\documentclass[11pt]{article}
\usepackage{amsmath,amssymb,amsthm,hyperref}
\usepackage[margin=26mm]{geometry}
\newtheorem{theorem}{Theorem}\newtheorem{lemma}{Lemma}
\newcommand{\Q}{\mathbb Q}\newcommand{\Z}{\mathbb Z}
\title{A rational elliptic denominator-support block}
\author{}\date{17 November 2026}
\begin{document}\maketitle

\section{Points and the predicate}

Use $E:y^2=x^3-36x$, whose rational group has rank one, generator
$(-3,9)$ and torsion subgroup $E[2](\Q)$; these data are recorded
in the primary database entry \cite{LMFDB}. Put
$P=2(-3,-9)=(25/4,35/8)$. The square-abscissa points are exactly
$2E(\Q)\setminus\{O\}=\Z P\setminus\{O\}$.
Indeed the two-isogeny descent classes of the torsion points are
$1,-1,6,-6$, and that of $(-3,9)$ is $-3$; an odd multiple,
even after adding torsion, has class $-3,3,-2$ or $2$.
Among the torsion classes only $1$ is a square. The converse follows
from the duplication formula.

Thus one retained rational coordinate $q$ and the finite quadratic
generator $r$ give a point by
\[
                 r^2=q^4-36,\qquad x=q^2,\quad y=qr.
\]
For two such points $H=[h]P$, $U=[j]P$, let their retained
coordinates be $q_H,q_U$. Define
\[
 u=\frac1{q_H},\quad C=\frac1{q_U^2},\quad
 D=\frac{C^2}{1+uC^2},\quad W=1+\frac{D^2}{u},\quad
 b=\frac{uW^2}{(1+u)^2},\quad\eta=7b.
\]
Guard all displayed denominators and $W$. Impose the intrinsic pole
equation with two retained parameters $v,s$ and one finite root $Y$:
\[
 Y^2=\eta(\eta+v^2)^2s^4
          -\eta(\eta-v^2)^2(1-\eta)^2.
\]

\begin{lemma}
Every rational solution has the following property at every prime:
if the denominator of $q_H$ has positive odd valuation, then $U$
reduces to the identity in the formal group.
Conversely, if $U=[n]H$ for an integer $n\ne0$, then $b$ has
square reduced denominator at every prime.
\end{lemma}
\begin{proof}
Let $\nu=v_p(u)$ be odd and $c=v_p(C)$. Distinct parity gives
\[
                    v_p(D)=\max(2c,-\nu).
\]
Suppose $\nu>0$ and $U$ is not formal. Then $v_p(q_U)\ge0$,
so $c\le0$ and $d_0:=v_p(D)\le0$. Therefore
$v_p(W)=2d_0-\nu$ and
\[
                     v_p(b)=4d_0-\nu<0
\]
is odd, contradicting intrinsic-pole soundness.

For the converse, first suppose $\nu>0$ is odd. Then $H$ is
formal, with $v_p(\zeta(H))=\nu$, where $\zeta=-x/y$.
The integral formal group gives
$v_p(\zeta([n]H))\ge\nu$, including at two and three.
Thus $c\ge2\nu$, $v_p(D)\ge4\nu$, and $W$ is a local unit.
It follows that $v_p(b)=\nu>0$. If $\nu<0$ is odd, the displayed
maximum instead gives $v_p(D^2/u)\ge-3\nu>0$, so $W$ is a unit
and $v_p(b)=-\nu>0$. Finally, even valuations of $u$ give even
valuations of $u/(1+u)^2$, and multiplication by $W^2$ preserves
parity. Hence no negative odd valuation of $b$ occurs.
\end{proof}

\section{Prepared completeness, uniformly in the quotient}

\begin{theorem}
For every fixed finite index cutoff $B$, there are infinitely many
points $H=[h]P$ with $|h|>B$ such that the displayed pole block
has rational witnesses for every $U=[n]H$, $n\in\Z\setminus\{0\}$.
Every retained whole-algebra guard can be satisfied.
\end{theorem}
\begin{proof}
Choose the positive real square root $q_H$ and require
\[
              \frac1{2\cdot6^4}<u<\frac2{6^4}.
\]
For every nonzero doubled rational point $U$, its positive square
abscissa satisfies $x(U)>6$, so $0<C<1/6$. Hence
$0<D<C^2<1/36$ and
\[
 0<b\le u\left(1+\frac1{6^4u}\right)^2
       =u+\frac2{6^4}+\frac1{6^8u}
       <\frac6{6^4}=\frac1{216}<\frac17.
\]
This real bound is uniform in the integer $n$.

Simultaneously require $u\in\Q_2^{\times2}$ with even valuation
at least four, and require $H$ sufficiently formal at seven.
These conditions are available in the cyclic group. To see the local
square-class freedom explicitly, write $H=2R$ with $R$ a multiple
of $(-3,-9)$. The duplication square root is
\[
 q_H=\frac{x(R)^2+36}{2y(R)},\qquad
 u=\frac{2y(R)}{x(R)^2+36}\sim-2\zeta(R)
\]
near the identity. A sufficiently deep fixed congruence on the
multiple of $R$ prescribes this $2$-adic square class and valuation,
as well as the required depth at seven. Any such congruence progression
is dense in its real connected cyclic closure; one can impose the
displayed positive interval and avoid finitely many indices.

For $U=[n]H$, the preceding lemma makes $W$ a unit at these prepared
places. Thus $b$ is a $2$-adic square of valuation at least four and
is seven-integral. Together with the square denominator and the real
bound, $\eta=7b$ satisfies the prepared intrinsic-pole theorem.
That theorem supplies Zariski-dense witness pairs $(v,s)$.

The values $u,C,D,W,b$ depend only on the retained coordinates
$q_H,q_U$, not on the finite ordinate signs. They are positive and
their denominators nonzero at every point-algebra character of these
prepared tuples. Pole-witness density avoids the remaining radicand
norm and separator guards.
\end{proof}

\section{Index divisibility and scope}

Under the simple-primitive-denominator theorem of \cite{Primitive},
weak rational $abc$ at exponent $16/15$ gives an index threshold
$n_0$. A finite guard can force $|h|>B\ge n_0$ without an additional
coordinate:
\[
 \prod_{m=1}^{B}\bigl(\operatorname{den}(x([m]P))q_H^2
                    -\operatorname{num}(x([m]P))\bigr)\ne0.
\]
The support implication then forces $h\mid j$. Its prepared
completeness supplies every nonzero integer quotient $j/h$.
The two point quadratics and one pole quadratic have finite rank
eight; one separator coordinate suffices for their guarded encoder.

This does not identify an external rational variable $X$ with $j/h$.
Congruences modulo the odd denominator support of $H$ leave the integer
quotient unbounded. A second source must establish the input link;
no nine-variable defining polynomial follows from this block alone.

\begin{thebibliography}{9}
\bibitem{LMFDB}The LMFDB Collaboration, elliptic curve $576.c3$,
equation $y^2=x^3-36x$, Mordell--Weil group and generators,
\url{https://www.lmfdb.org/EllipticCurve/Q/576/c/3}, accessed
17 November 2026.
\bibitem{Primitive}\emph{Simple primitive denominator primes under
weak $abc$}, companion checkpoint manuscript,
\nolinkurl{checkpoints/general_scalar_v56/simple_primitive_denominators_v59.tex}.
\end{thebibliography}
\end{document}
